%%% Please fill in basic information on your thesis, which will be automatically
%%% inserted at the right places. You need to replace \xxx{...} by real data.

% Type of your thesis:
%	"bc" for Bachelor's
%	"mgr" for Master's
%	"phd" for PhD
%	"rig" for rigorosum
\def\ThesisType{bc}

% Language of your study programme:
%	"cs" for Czech
%	"en" for English
\def\StudyLanguage{cs}

% Thesis title in English (exactly as in the official assignment)
% (Note: \xxx is a "ToDo label" which makes the unfilled visible. Remove it.)
\def\ThesisTitle{\xxx{Measuring Morphological Complexity of~Languages Using Broadly Multilingual Data}}

% Author of the thesis (you)
\def\ThesisAuthor{\xxx{Antonín Forst}}

% Year when the thesis is submitted
\def\YearSubmitted{\xxx{2026}}

% Name of the department or institute, where the work was officially assigned
% (according to the Organizational Structure of MFF UK in English,
% see https://www.mff.cuni.cz/en/faculty/organizational-structure,
% or a full name of a department outside MFF)
\def\Department{\xxx{Institute of Formal and Applied Linguistics}}

% Is it a department (katedra), or an institute (ústav)?
\def\DeptType{\xxx{Institute}}

% Thesis supervisor: name, surname and titles
\def\Supervisor{\xxx{prof. Ing. Zdeněk Žabokrtský, Ph.D.}}

% Supervisor's department (again according to Organizational structure of MFF)
\def\SupervisorsDepartment{\xxx{Institute of Formal and Applied Linguistics}}

% Study programme (does not apply to rigorosum theses)
\def\StudyProgramme{\xxx{Artifical inteligence}}

% An optional dedication: you can thank whomever you wish (your supervisor,
% consultant, who provided you with tea and pizza, etc.)
\def\Dedication{

At this place, I would especially like to thank my supervisor, Professor Zdeněk Žabokrtský, for his patient work with me. His advice and guidance helped me not only with the factual and professional aspects of the thesis, but also more generally with the process of preparing a larger coherent piece of work.

I would also like to thank Mgr. Vojtěch John for providing the morphological segmentation of the data essential for this thesis.

Last but not least, I would like to thank my family, who continue to support me in my seemingly never-ending studies

Thank you.
}


% Abstract (recommended length around 80-200 words; this is not a copy of your thesis assignment!)
\def\Abstract{%
\xxx{
The thesis investigates morphological complexity of languages using broadly multilingual data. It obtains and preprocesses corpus data for 443 languages from multiple sources. It applies various complexity metrics to a 175 of them. The metrics are based on corpus data. The thesis further proposes metrics for morphological segmentation and applies them on segmented data for 40 languages. It required to process gigabytes of raw linguistic data of different quality, different scripts, and other properties. The results are analyzed and compared with each other, applied to other linguistic categories and showing various patterns and implications on multiple different presented graphs.
}
}

% 3 to 5 keywords (recommended) separated by \sep
% Keywords are useful for indexing and searching for the theses by topic.
\def\ThesisKeywords{%
\xxx{
morphological complexity \sep
complexity metrics \sep
broadly multilingual data \sep
corpus linguisitics \sep
morphological segmentation
}
}

% If any of your metadata strings contains TeX macros, you need to provide
% a plain-text version for use in XMP metadata embedded in the output PDF file.
% If you are not sure, check the generated thesis.xmpdata file.
\def\ThesisAuthorXMP{\ThesisAuthor}
\def\ThesisTitleXMP{\ThesisTitle}
\def\ThesisKeywordsXMP{\ThesisKeywords}
\def\AbstractXMP{\Abstract}

% If your abstracts are long and do not fit in the infopage, you can make the
% fonts a bit smaller by this setting. (Also, you should try to compress your abstract more.)
\def\InfoPageFont{}
%\def\InfoPageFont{\small}  % uncomment to decrease font size

% If you are studing in a Czech programme, you also need to provide metadata in Czech:
% (in English programmes, this is not used anywhere)

\def\ThesisTitleCS{\xxx{Měření morfologické komplexity jazyků na základě široce multilingválních dat}}
\def\DepartmentCS{\xxx{Ústav formální a aplikované lingvistiky}}
\def\DeptTypeCS{\xxx{Ústav}}
\def\SupervisorsDepartmentCS{\xxx{Ústav formální a aplikované lingvistiky}}
\def\StudyProgrammeCS{\xxx{Umělá inteligence}}

\def\ThesisKeywordsCS{%
\xxx{
morfologická komplexita \sep
metriky komplexity \sep
široce multilingvální data \sep
korpusová lingvistika \sep
morfologická segmentace
}
}

\def\AbstractCS{%
\xxx{Tato práce zkoumá morfologickou komplexitu jazyků s využitím široce vícejazyčných dat. Získává a předzpracovává korpusová data pro 443 jazyků z několika zdrojů. Na 175 z nich aplikuje různé metriky komplexity založené na korpusových datech. Práce dále navrhuje metriky pro morfologickou segmentaci a aplikuje je na segmentovaná data pro 40 jazyků. Pro tyto účely bylo potřeba zpracovat gigabyty dat různé kvality, písem a dalších vlastností. Výsledky jsou analyzovány a vzájemně porovnány, aplikovány na další lingvistické kategorie a jsou zdůrazněny různé vzory a jejich důsledky na přiložených grafech.}
}
